\documentclass[12pt]{article}
\usepackage[spanish]{babel}
\usepackage[T1]{fontenc}
\usepackage[utf8]{inputenc}
\usepackage{amsmath,amssymb}
\usepackage{amsthm}
\usepackage{geometry}
\geometry{margin=2.5cm}
\usepackage{comment}
\renewcommand{\familydefault}{\sfdefault}
\usepackage{mathrsfs}

% Entornos de teorema (mismo estilo que CI.pdf)
\newtheorem{definicion}{Definición}[section]
\newtheorem{proposicion}{Proposición}[section]
\newtheorem{teorema}{Teorema}[section]
\newtheorem{corolario}{Corolario}[section]
\newtheorem{observacion}{Observación}[section]
\newtheorem{ejemplo}{Ejemplo}[section]

% Entornos nuevos para este documento
\newtheorem{conjetura}{Conjetura}[section]
\newtheorem{problema}{Problema abierto}[section]

% Figuras y tablas
\usepackage{graphicx}
\usepackage{float}
\usepackage{subcaption}
\usepackage{booktabs,tabularx,array}
\usepackage{longtable}
\usepackage{xcolor}
\usepackage{colortbl}

% Listas
\usepackage{enumitem}

% Títulos
\usepackage{titlesec}

% Bibliografía
\usepackage[authoryear]{natbib}
\bibliographystyle{plainnat}

% Hipervínculos
\usepackage[hidelinks]{hyperref}

% Colores para la tabla de orientación
\definecolor{colorDef}{HTML}{E6F1FB}
\definecolor{colorProp}{HTML}{EEEDFE}
\definecolor{colorConj}{HTML}{FBEAF0}
\definecolor{colorProb}{HTML}{F1EFE8}
\definecolor{colorExp}{HTML}{FAECE7}
\definecolor{trivBaja}{HTML}{EAF3DE}
\definecolor{trivMedia}{HTML}{FAEEDA}
\definecolor{trivAlta}{HTML}{FCEBEB}
\definecolor{trivAbierto}{HTML}{E1F5EE}

\title{Control Informacional\\
\large Proyecto de Investigación II --- Esquema de trabajo}

\author{
\textbf{Alejandro Cruz López} \\
Universidad Autónoma Metropolitana \\
Unidad Iztapalapa
}

\date{
\vspace{1em}
\textbf{Asesores} \\[0.3em]
Dr.\ Asael Fabián Martínez Martínez \\
Dr.\ Joaquín Delgado Fernández \\
\vspace{1em}
2026
}

\begin{document}
\maketitle

\begin{abstract}
Este documento presenta el esquema de trabajo del Proyecto de Investigación~II,
cuyo eje es el estudio de \emph{dinámicas restringidas sobre el símplex de
distribuciones de probabilidad discretas}. El proyecto se organiza en tres
problemas centrales: alcanzabilidad restringida, control óptimo informacional y
defecto de ordenamiento. Para cada problema se enuncian las definiciones,
proposiciones, conjeturas y problemas abiertos pertinentes, acompañados de un
esbozo breve de la estrategia de demostración. El documento concluye con una
tabla de orientación que resume los 23~elementos del esquema junto con su nivel
de trivialidad y rol en el proyecto.
\end{abstract}

\tableofcontents
\newpage

%====================================================================
\section{Preliminares y notación heredada de CI}
%====================================================================

Se conserva íntegramente la notación del Proyecto~I. En particular:
\begin{itemize}[nosep]
  \item $\Delta_n$ es el símplex de probabilidad y $\Delta_n^\circ$ su interior.
  \item $T_Y$ es el operador informacional global con $Y\in C_{\mathrm{glob}}$.
  \item $L_K(X)=XK$ es el operador de Markov asociado a la matriz estocástica~$K$.
  \item $B_L(X)_i = X_iL_i/\langle X,L\rangle$ es el operador bayesiano con
        evidencia $L\in\mathbb{R}_{>0}^n$.
  \item $D_{\mathrm{KL}}(P\|X)=\sum_i P_i\log(P_i/X_i)$ es la divergencia de
        Kullback--Leibler.
  \item $g_X(u,v)=\sum_i u_iv_i/X_i$ es la métrica de Fisher en $T_X\Delta_n^\circ$.
\end{itemize}

La diferencia central respecto al Proyecto~I es que ahora se imponen
\emph{restricciones} sobre las acciones admisibles. Sin ellas, tanto $B_L$ como
$L_K$ pueden llevar cualquier estado a cualquier objetivo en un solo paso, lo que
hace trivial todo problema de control o alcanzabilidad.

%====================================================================
\section{Problema 1 --- Alcanzabilidad restringida}
%====================================================================

\noindent\textbf{Pregunta rectora.} \textit{¿Es $P$ alcanzable desde $X_0$ en
$t$ pasos bajo acciones admisibles $({\mathcal K}_\varepsilon, {\mathcal U}_R)$?}

\bigskip

\begin{definicion}[Acciones admisibles]
\label{def:acciones}
Se definen los conjuntos de acciones admisibles como
\[
\mathcal{K}_\varepsilon
:=
\bigl\{K\in\mathcal{S}_n : K_{ij}\geq\varepsilon>0\text{ para todo }i,j\bigr\}
\]
y
\[
\mathcal{U}_R
:=
\bigl\{v\in\mathbb{R}^n : \|v\|_\infty\leq R\bigr\},
\]
donde $\varepsilon>0$ y $R>0$ son parámetros fijos. Las acciones de Markov
admisibles son $K\in\mathcal{K}_\varepsilon$ y las bayesianas son evidencias $L$
con $\log L\in\mathcal{U}_R$.
\end{definicion}

\begin{observacion}
La condición $K_{ij}\geq\varepsilon$ impide matrices de rango~$1$ (que llevarían
cualquier $X$ al objetivo exacto en un paso) y garantiza que toda
$K\in\mathcal{K}_\varepsilon$ sea primitiva. La condición $\|\log L\|_\infty\leq R$
impide la evidencia dirigida exacta $L_i=P_i/X_i$. Ambas restricciones son las
que hacen no trivial el problema de alcanzabilidad.
\end{observacion}

\begin{definicion}[Costo informacional de una acción]
\label{def:costo}
Para $X\in\Delta_n^\circ$, se define el costo informacional de una acción como
\[
c_M(K,X):=D_{\mathrm{KL}}(XK\|X),
\qquad
c_B(L,X):=D_{\mathrm{KL}}(B_L(X)\|X).
\]
Ambos costos miden cuánto desplaza la acción al estado actual en términos
informacionales, son no negativos, y se anulan cuando la acción es la identidad.
\end{definicion}

\begin{observacion}
La elección de $D_{\mathrm{KL}}$ como costo hace comparables las acciones de
Markov y Bayes en una escala común. Esta elección es la que permite escribir
el funcional de costo del Problema~2 con una penalización uniforme.
\end{observacion}

\begin{definicion}[Conjunto alcanzable restringido]
\label{def:alcanzable}
Dado $X_0\in\Delta_n^\circ$ y un horizonte $t\geq 1$, el \emph{conjunto
alcanzable restringido en $t$ pasos} es
\[
\mathcal{R}_t(X_0)
:=
\bigl\{X_t\in\Delta_n^\circ :
X_{s+1}\in\{L_{K_s}(X_s),\, B_{L_s}(X_s)\},\;
K_s\in\mathcal{K}_\varepsilon,\;
\log L_s\in\mathcal{U}_R
\bigr\}.
\]
\end{definicion}

\begin{definicion}[Conjunto alcanzable cuantizado]
\label{def:alcanzable-cuant}
Dado $M\geq 1$, el \emph{retículo de resolución $M$} es
\[
\Delta_n^{(M)}
:=
\Bigl\{X\in\Delta_n : MX_i\in\mathbb{Z}_{\geq 0}\text{ para todo }i\Bigr\}.
\]
El \emph{conjunto alcanzable cuantizado en $t$ pasos} es
\[
\mathcal{R}_t^{(M)}(X_0)
:=
\mathcal{R}_t(X_0)\cap\Delta_n^{(M)},
\]
donde además se exige que todos los estados intermedios pertenezcan a
$\Delta_n^{(M)}$.
\end{definicion}

\begin{proposicion}[Inalcanzabilidad en tiempo finito]
\label{prop:inalcanzable}
Bajo las restricciones $(\mathcal{K}_\varepsilon,\mathcal{U}_R)$, para todo
$X_0\in\Delta_n^\circ$ y todo $t$ finito existen puntos $P\in\Delta_n^\circ$
tales que $P\notin\mathcal{R}_t(X_0)$.
\end{proposicion}

\begin{proof}[Esbozo]
Cada acción admisible de Markov mueve el estado en una dirección acotada:
$\|L_K(X)-X\|_1\leq 2(1-n\varepsilon)$ para $K\in\mathcal{K}_\varepsilon$.
Cada acción bayesiana admisible satisface $\|B_L(X)-X\|_1\leq C(R)$ para una
constante $C(R)$ que depende de $R$ pero no de $X$. Por tanto, en $t$ pasos la
distancia en norma~$\ell^1$ entre $X_t$ y $X_0$ está uniformemente acotada.
Escogiendo $P$ a distancia mayor que esa cota se concluye la inalcanzabilidad.
\end{proof}

\begin{problema}[Caracterización del conjunto alcanzable]
\label{prob:caracterizacion}
Determinar condiciones necesarias y suficientes sobre $(X_0, P,
\mathcal{K}_\varepsilon, \mathcal{U}_R, t)$ para que $P\in\mathcal{R}_t(X_0)$.
En particular, caracterizar la geometría de $\mathcal{R}_t(X_0)$ como
subconjunto de $\Delta_n^\circ$.
\end{problema}

\begin{conjetura}[Pérdida de alcanzabilidad por cuantización]
\label{conj:cuantizacion}
Para $M$ suficientemente grande y bajo las restricciones
$(\mathcal{K}_\varepsilon,\mathcal{U}_R)$, la pérdida de alcanzabilidad debida a
la cuantización es de orden $O(1/M)$ en distancia~$D_{\mathrm{KL}}$. Más
precisamente, existe $C>0$ tal que para todo $P\in\mathcal{R}_t(X_0)$,
\[
\inf_{Q\in\mathcal{R}_t^{(M)}(X_0)}D_{\mathrm{KL}}(P\|Q)\leq \frac{C}{M}.
\]
\end{conjetura}

%====================================================================
\section{Problema 2 --- Control óptimo informacional}
%====================================================================

\noindent\textbf{Pregunta rectora.} \textit{¿Cuándo conviene aplicar Markov,
cuándo Bayes y cuándo alternar?}

\bigskip

\begin{definicion}[Política admisible de horizonte $T$]
\label{def:politica}
Una \emph{política admisible de horizonte $T$} es una sucesión de decisiones
$\{a_t\}_{t=0}^{T-1}$ con $a_t\in\{M,B\}$, acompañada de parámetros
$K_t\in\mathcal{K}_\varepsilon$ si $a_t=M$, o $\log L_t\in\mathcal{U}_R$ si
$a_t=B$. El estado evoluciona según
\[
X_{t+1}
=
\begin{cases}
L_{K_t}(X_t) & \text{si }a_t=M, \\
B_{L_t}(X_t) & \text{si }a_t=B.
\end{cases}
\]
\end{definicion}

\begin{definicion}[Funcional de costo]
\label{def:costo-total}
Dado un objetivo $P\in\Delta_n^\circ$ y un parámetro de regularización
$\lambda>0$, el \emph{costo de una política admisible} es
\[
J_T(X_0;\{a_t\})
:=
D_{\mathrm{KL}}(P\|X_T)
+
\lambda\sum_{t=0}^{T-1}c_{a_t}(a_t,X_t),
\]
donde $c_{a_t}$ es el costo de acción de la Definición~\ref{def:costo}.
\end{definicion}

\begin{definicion}[Función de valor óptimo]
\label{def:valor}
La \emph{función de valor óptimo} con $\tau$ pasos restantes es
\[
V_\tau(X)
:=
\inf_{\{a_t\}_{t=0}^{\tau-1}\text{ admisible}} J_\tau(X;\{a_t\}),
\qquad
V_0(X):=D_{\mathrm{KL}}(P\|X).
\]
\end{definicion}

\begin{proposicion}[Ecuación de Bellman informacional restringida]
\label{prop:bellman}
La función de valor satisface la recursión hacia atrás
\[
V_\tau(X)
=
D_{\mathrm{KL}}(P\|X)
+
\min\!\left(
\inf_{K\in\mathcal{K}_\varepsilon}
\bigl[\lambda\,c_M(K,X)+V_{\tau-1}(XK)\bigr],\;
\inf_{\log L\in\mathcal{U}_R}
\bigl[\lambda\,c_B(L,X)+V_{\tau-1}(B_L(X))\bigr]
\right).
\]
Bajo las restricciones $(\mathcal{K}_\varepsilon,\mathcal{U}_R)$, el ínfimo no es
idénticamente cero para $\tau$ finito.
\end{proposicion}

\begin{proof}[Esbozo]
La recursión se obtiene por el principio de optimalidad de Bellman: el costo
óptimo desde $X$ en $\tau$ pasos es el costo del primer paso más el costo óptimo
desde el estado resultante en $\tau-1$ pasos. La no trivialidad del ínfimo se
sigue de la Proposición~\ref{prop:inalcanzable}: si $P\notin\mathcal{R}_1(X)$,
entonces $V_1(X)>0$ para toda acción admisible.
\end{proof}

\begin{conjetura}[Política de tipo \emph{bang-bang} para $n=2$]
\label{conj:bangbang}
Para $n=2$ y restricciones simétricas
($\varepsilon$ y $R$ iguales para ambas acciones), existe una partición de
$\Delta_2^\circ$ en dos regiones $\mathcal{M}$ y $\mathcal{B}$ tales que la
política óptima es $a^*=M$ en $\mathcal{M}$ y $a^*=B$ en $\mathcal{B}$. En
particular, la política óptima es de acción pura en cada región y no requiere
alternancia.
\end{conjetura}

\begin{conjetura}[Ventaja de la alternancia para $n\geq 3$]
\label{conj:alternancia}
Para $n\geq 3$, existen configuraciones $(X_0, P, \mathcal{K}_\varepsilon,
\mathcal{U}_R)$ para las cuales
\[
J_T^*(X_0)<\min\!\bigl(J_T^M(X_0),\,J_T^B(X_0)\bigr),
\]
donde $J_T^*$ es el costo de la política óptima y $J_T^M$, $J_T^B$ son los
costos de las políticas puras (solo Markov y solo Bayes respectivamente). Es
decir, la alternancia de operadores produce un costo estrictamente menor que
cualquier política pura.
\end{conjetura}

\begin{conjetura}[Límite continuo igual al flujo Fisher--KL]
\label{conj:limite-continuo}
Para $T$ fijo y $\Delta t=T/N$ con costo de acción $c_M(K,X)$, en el límite
$N\to\infty$ la trayectoria de la política óptima converge a la solución de la
ecuación diferencial
\[
\dot{X}=P-X,
\]
que es el flujo de gradiente Fisher--KL de $D_{\mathrm{KL}}(P\|\cdot)$
establecido en CI~(Proyecto~I).
\end{conjetura}

%====================================================================
\section{Problema 3 --- Defecto de ordenamiento}
%====================================================================

\noindent\textbf{Pregunta rectora.} \textit{¿Cuánto cambia la trayectoria según
el orden en que se aplican Markov y Bayes?}

\bigskip

\begin{definicion}[Defecto de ordenamiento]
\label{def:defecto}
Sea $K\in\mathcal{K}_\varepsilon$ y sea $L$ con $\log L\in\mathcal{U}_R$. El
\emph{defecto de ordenamiento} en el estado $X\in\Delta_n^\circ$ es el vector
\[
\mathcal{D}_{K,L}(X)
:=
B_L(XK)-B_L(X)K
\in\mathbb{R}^n,
\]
que mide la diferencia entre aplicar primero Markov y luego Bayes, versus el
orden inverso.
\end{definicion}

\begin{observacion}
El defecto $\mathcal{D}_{K,L}(X)$ generaliza y reemplaza el conmutador
$\mathcal{C}_{P,L}(X)$ de CI~(Proyecto~I). La diferencia es de convención: aquí
se usa $K$ para la matriz de Markov (reservando $P$ para el objetivo) y se
escribe la diferencia en el orden que resulta más natural para el problema de
control.
\end{observacion}

\begin{proposicion}[Tangencia al símplex del defecto]
\label{prop:tangencia}
Para todo $X\in\Delta_n^\circ$, toda $K\in\mathcal{K}_\varepsilon$ y toda $L$
con $\log L\in\mathcal{U}_R$, se cumple
\[
\sum_{i=1}^n \mathcal{D}_{K,L}(X)_i=0.
\]
Por tanto,
\[
\mathcal{D}_{K,L}(X)\in T_X\Delta_n^\circ.
\]
\end{proposicion}

\begin{proof}[Esbozo]
Ambos términos en la definición del defecto son distribuciones de probabilidad:
$B_L(XK)\in\Delta_n^\circ$ y $B_L(X)K\in\Delta_n^\circ$. Por tanto, cada
uno tiene suma de coordenadas igual a~$1$. La diferencia tiene suma cero, lo
que es precisamente la condición de pertenencia al espacio tangente
$T_X\Delta_n^\circ=\{u\in\mathbb{R}^n:\sum_i u_i=0\}$.
\end{proof}

\begin{problema}[Caracterización algebraica de $\mathcal{D}_{K,L}\equiv 0$]
\label{prob:conmutacion}
Determinar condiciones necesarias y suficientes sobre $K$ y $L$ para que
\[
\mathcal{D}_{K,L}(X)=0
\quad\text{para todo }X\in\Delta_n^\circ.
\]
\end{problema}

\begin{observacion}
El Proyecto~I contenía la condición $L\propto\pi$ (con $\pi$ estacionaria de~$K$)
como posible respuesta. Sin embargo, esa condición es insuficiente en general con
la convención de vectores fila. El Problema~\ref{prob:conmutacion} queda abierto
y su resolución es necesaria antes de usar cualquier resultado de conmutación
en el control óptimo.
\end{observacion}

\begin{proposicion}[Subvariedad de conmutación]
\label{prop:subvariedad}
Bajo condiciones genéricas sobre $K$ y $L$, el conjunto
\[
\mathcal{M}_{K,L}
:=
\bigl\{X\in\Delta_n^\circ:\mathcal{D}_{K,L}(X)=0\bigr\}
\]
es una subvariedad de $\Delta_n^\circ$ de dimensión estrictamente menor que
$n-1$.
\end{proposicion}

\begin{proof}[Esbozo]
El defecto $\mathcal{D}_{K,L}:\Delta_n^\circ\to T\Delta_n^\circ$ es una función
suave. Por el teorema de la función implícita, si la diferencial
$D_X\mathcal{D}_{K,L}$ tiene rango máximo en algún punto de
$\mathcal{M}_{K,L}$, entonces $\mathcal{M}_{K,L}$ es localmente una subvariedad.
La dimensión es estrictamente menor que $n-1$ porque $\mathcal{D}_{K,L}=0$ impone
$n-1$ ecuaciones independientes en un espacio de dimensión $n-1$ (el símplex).
La genericidad de las condiciones debe verificarse por separado.
\end{proof}

\begin{conjetura}[Cota de suboptimalidad por orden incorrecto]
\label{conj:cota}
Existe una constante $C=C(\varepsilon,R)>0$ tal que para toda $K\in
\mathcal{K}_\varepsilon$, toda $L$ con $\log L\in\mathcal{U}_R$ y todo
$X\in\Delta_n^\circ$,
\[
\bigl|J(X;K,L)-J(X;L,K)\bigr|
\leq
C\cdot\|\mathcal{D}_{K,L}(X)\|,
\]
donde $J(X;K,L):=D_{\mathrm{KL}}(P\|B_L(XK))$ y
$J(X;L,K):=D_{\mathrm{KL}}(P\|B_L(X)K)$ son los costos de dos pasos en cada
orden.
En particular, si $\mathcal{D}_{K,L}(X)=0$, el orden de aplicación es
irrelevante para el costo.
\end{conjetura}

\begin{observacion}[El corazón del proyecto]
La Conjetura~\ref{conj:cota} es el resultado central que unifica los tres
problemas:
\begin{itemize}[nosep]
  \item \textbf{Problema~1:} los dos órdenes producen estados distintos,
        cuya separación está medida por $\mathcal{D}_{K,L}(X)$.
  \item \textbf{Problema~2:} la diferencia de costos entre el orden óptimo y
        el subóptimo es precisamente $|J(X;K,L)-J(X;L,K)|$.
  \item \textbf{Problema~3:} el defecto $\mathcal{D}_{K,L}$ es el objeto que
        acota esa diferencia.
\end{itemize}
Una demostración con constante $C$ no óptima, o bajo condiciones adicionales
sobre $K$ y $L$, ya constituiría un resultado publicable.
\end{observacion}

%====================================================================
\section{Tabla de orientación}
%====================================================================

La tabla siguiente resume los 23~elementos del esquema con su tipo, contenido
abreviado, rol en el proyecto y nivel de trivialidad. Su propósito es servir
como mapa de navegación durante la redacción formal del documento.

\medskip

\noindent\textbf{Leyenda de tipos:}
\textbf{Def}~=~Definición,\;
\textbf{Prop}~=~Proposición,\;
\textbf{Conj}~=~Conjetura,\;
\textbf{Prob}~=~Problema abierto,\;
\textbf{Exp}~=~Experimento.

\medskip

\noindent\textbf{Leyenda de trivialidad:}
\colorbox{trivBaja}{\strut~Baja~} formalización necesaria, no sorprendente;\;
\colorbox{trivMedia}{\strut~Media~} requiere demostración;\;
\colorbox{trivAlta}{\strut~Alta~} difícil;\;
\colorbox{trivAbierto}{\strut~Abierto~} sin respuesta conocida.

\bigskip

\begingroup
\small
\setlength{\tabcolsep}{4pt}
\renewcommand{\arraystretch}{1.25}

\noindent\begin{longtable}{@{}p{0.65cm}p{0.65cm}p{3.5cm}p{4.6cm}p{3.6cm}p{1.5cm}@{}}
\toprule
\textbf{Tipo} & \textbf{ID} & \textbf{Título} & \textbf{Contenido} &
\textbf{Rol en el proyecto} & \textbf{Trivial.} \\
\midrule
\endfirsthead
\multicolumn{6}{c}{\textit{(continuación)}} \\
\toprule
\textbf{Tipo} & \textbf{ID} & \textbf{Título} & \textbf{Contenido} &
\textbf{Rol en el proyecto} & \textbf{Trivial.} \\
\midrule
\endhead
\bottomrule
\endlastfoot

\multicolumn{6}{@{}l}{\textbf{Problema 1 --- Alcanzabilidad restringida}} \\[2pt]
\rowcolor{colorDef}
Def & 1.1 & Acciones admisibles $({\mathcal K}_\varepsilon,{\mathcal U}_R)$ &
$K_{ij}\geq\varepsilon$; $\|\log L\|_\infty\leq R$ &
Sin esto todo colapsa a un paso &
\cellcolor{trivBaja} Baja \\
\rowcolor{colorDef}
Def & 1.2 & Costo informacional de acción &
$c_M=D_{\mathrm{KL}}(XK\|X)$; $c_B=D_{\mathrm{KL}}(B_L(X)\|X)$ &
Hace comparables Markov y Bayes en escala KL &
\cellcolor{trivBaja} Baja \\
\rowcolor{colorDef}
Def & 1.3 & Conjunto alcanzable $\mathcal{R}_t(X_0)$ &
Estados alcanzables en $t$ pasos bajo $({\mathcal K}_\varepsilon,{\mathcal U}_R)$ &
Objeto central del Problema~1 &
\cellcolor{trivBaja} Baja \\
\rowcolor{colorDef}
Def & 1.4 & Conjunto alcanzable cuantizado $\mathcal{R}_t^{(M)}(X_0)$ &
$\mathcal{R}_t\cap\Delta_n^{(M)}$; estados intermedios en el retículo &
Restricción sobre estados, no solo acciones &
\cellcolor{trivBaja} Baja \\
\rowcolor{colorProp}
Prop & 1.1 & Inalcanzabilidad en tiempo finito &
Bajo $({\mathcal K}_\varepsilon,{\mathcal U}_R)$, existen $P\notin\mathcal{R}_t(X_0)$ &
Justifica la no trivialidad del marco restringido &
\cellcolor{trivMedia} Media \\
\rowcolor{colorProb}
Prob & 1.2 & Caracterización de $\mathcal{R}_t(X_0)$ &
CNS para $P\in\mathcal{R}_t(X_0)$ en función de $X_0,P,\varepsilon,R,t$ &
Problema abierto central; atacable con punto fijo &
\cellcolor{trivAbierto} Abierto \\
\rowcolor{colorConj}
Conj & 1.1 & Pérdida por cuantización $O(1/M)$ &
$\inf_{Q\in\mathcal{R}_t^{(M)}}D_{\mathrm{KL}}(P\|Q)\leq C/M$ &
Verificable numéricamente antes del ataque analítico &
\cellcolor{trivMedia} Media \\
\rowcolor{colorExp}
Exp & 1.A & Mapa de $\mathcal{R}_1$ y $\mathcal{R}_2$ para $n=3$ &
Trazar $\mathcal{R}_t(X_0)$ para $\varepsilon=0.05,0.1$ y $R=0.5,1,2$ &
Diseñado para refutar la Prop.~1.1 o acotar la constante &
--- \\[4pt]

\multicolumn{6}{@{}l}{\textbf{Problema 2 --- Control óptimo informacional}} \\[2pt]
\rowcolor{colorDef}
Def & 2.1 & Política admisible de horizonte $T$ &
$\{a_t\}$ con $a_t\in\{M,B\}$ y parámetros en $({\mathcal K}_\varepsilon,{\mathcal U}_R)$ &
Formaliza el espacio de decisiones &
\cellcolor{trivBaja} Baja \\
\rowcolor{colorDef}
Def & 2.2 & Funcional de costo $J_T$ &
$D_{\mathrm{KL}}(P\|X_T)+\lambda\sum_t c_{a_t}(a_t,X_t)$; indexación corregida &
Objeto de minimización del Problema~2 &
\cellcolor{trivBaja} Baja \\
\rowcolor{colorDef}
Def & 2.3 & Función de valor óptimo $V_T(X)$ &
$\inf_{\{a_t\}} J_T$; definida recursivamente hacia atrás &
Permite escribir la ecuación de Bellman &
\cellcolor{trivBaja} Baja \\
\rowcolor{colorProp}
Prop & 2.1 & Ecuación de Bellman restringida &
Recursión donde el ínfimo se toma sobre ${\mathcal K}_\varepsilon$ y ${\mathcal U}_R$, no sobre todos los operadores &
El ínfimo no es cero; distingue del caso irrestricto &
\cellcolor{trivMedia} Media \\
\rowcolor{colorConj}
Conj & 2.1 & Política \emph{bang-bang} para $n=2$ &
Partición de $\Delta_2^\circ$ en regiones de acción pura (solo $M$ o solo $B$) &
Caso tratable analíticamente; punto de entrada &
\cellcolor{trivMedia} Media \\
\rowcolor{colorConj}
Conj & 2.2 & Ventaja de la alternancia para $n\geq 3$ &
Existen $(X_0,P,{\mathcal K}_\varepsilon,{\mathcal U}_R)$ donde la política alternada supera a toda política pura &
Resultado central del Problema~2 &
\cellcolor{trivAbierto} Abierto \\
\rowcolor{colorConj}
Conj & 2.3 & Límite continuo $=$ flujo Fisher--KL &
Para $\Delta t\to 0$ con costo $c_M$, la política óptima converge a $\dot{X}=P-X$ &
Conecta control discreto con geometría de CI (Proyecto~I) &
\cellcolor{trivMedia} Media \\
\rowcolor{colorExp}
Exp & 2.A & Programación dinámica para $n=2$ &
Resolver $V_T$ sobre grilla de $\Delta_2^\circ$; trazar frontera entre regiones &
Diseñado para refutar la Conj.~2.1 &
--- \\
\rowcolor{colorExp}
Exp & 2.B & Comparación de políticas para $n=3$ &
Comparar costos: solo-$M$, solo-$B$, alternada, óptima numérica vs.\ $\varepsilon$ y $R$ &
Diseñado para refutar la Conj.~2.2 &
--- \\[4pt]

\multicolumn{6}{@{}l}{\textbf{Problema 3 --- Defecto de ordenamiento}} \\[2pt]
\rowcolor{colorDef}
Def & 3.1 & Defecto de ordenamiento $\mathcal{D}_{K,L}(X)$ &
$B_L(XK)-B_L(X)K$; diferencia entre los dos órdenes &
Observable central del Problema~3 &
\cellcolor{trivBaja} Baja \\
\rowcolor{colorProp}
Prop & 3.1 & Tangencia al símplex del defecto &
$\sum_i\mathcal{D}_{K,L}(X)_i=0$; verificación con convención fila &
Corrige el gap pendiente de CI (Proyecto~I) &
\cellcolor{trivBaja} Baja \\
\rowcolor{colorProb}
Prob & 3.2 & CNS algebraica de $\mathcal{D}_{K,L}\equiv 0$ &
Condiciones sobre $K$ y $L$ para conmutación global &
Reemplaza la condición incorrecta $L\propto\pi$ de CI &
\cellcolor{trivAlta} Alta \\
\rowcolor{colorProp}
Prop & 3.2 & Subvariedad de conmutación $\mathcal{M}_{K,L}$ &
$\{X:\mathcal{D}_{K,L}(X)=0\}$ es subvariedad de dimensión $<n-1$ &
Caracterización geométrica; por teorema de la función implícita &
\cellcolor{trivMedia} Media \\
\rowcolor{colorExp}
Exp & 3.A & Mapa del defecto en la grilla del símplex &
Calcular $\|\mathcal{D}_{K,L}(X)\|$ en grilla para $n=3$; identificar máximos y ceros &
Diseñado para refutar la Prop.~3.2 o verificar la dimensión de $\mathcal{M}_{K,L}$ &
--- \\
\rowcolor{colorConj}
$\bigstar$ Conj & 3.3 & Cota de suboptimalidad por orden incorrecto &
$|J(X;K,L)-J(X;L,K)|\leq C\|\mathcal{D}_{K,L}(X)\|$ &
\textbf{Clímax del proyecto:} une los tres problemas &
\cellcolor{trivAbierto} Abierto \\

\bottomrule
\end{longtable}
\endgroup

\bigskip

\noindent La Conjetura~$\bigstar$~3.3 es el objetivo principal del proyecto.
Una demostración con constante~$C$ no óptima, o bajo condiciones adicionales
sobre $K$ y~$L$, ya constituiría un resultado publicable. Si se refuta
numéricamente, la conclusión también es valiosa: indicaría que la suboptimalidad
por orden incorrecto no está controlada uniformemente por el defecto, lo que
requiere un observable distinto.

\end{document}
